\documentclass[12pt]{article}

\usepackage[utf8]{inputenc}
\usepackage[T1]{fontenc}
\usepackage{amsmath, amssymb}
\usepackage{graphicx}
\usepackage{geometry}
\usepackage{booktabs}
\usepackage{enumitem}

\geometry{margin=2.5cm}
\begin{document}

\section*{Guía Metodológica y Técnica Integral}
\addcontentsline{toc}{section}{Guía Metodológica y Técnica Integral}
\begin{center}
\large Análisis de Mercado y Herramienta de Recolección de Datos\\
(Enfoque ARCUSUR)
\end{center}

\section{Presentación y Enunciado Director}

Esta guía metodológica y técnica se diseña para orientar la ejecución de un
Análisis de Mercado sobre el perfil, expectativas de formación y
empleabilidad de los titulados de la carrera de Ingeniería en Sistemas. Los
resultados obtenidos servirán como insumo empírico fundamental para
alimentar el Plan de Desarrollo Estratégico y respaldar el proceso de
acreditación ante el sistema ARCUSUR.

\textbf{Enunciado Director:} Desarrollar un \emph{pipeline} de análisis
estadístico descriptivo, exploratorio, de ingeniería de software y
experimentación por simulación de sistemas sobre un conjunto estructurado
de datos de titulados profesionales de ingeniería de sistemas, con el fin
de cuantificar tendencias de formación continua, identificar patrones
demográficos y de preferencia académica, y proveer métricas robustas que
justifiquen la actualización curricular y la gestión de alianzas
estratégicas institucionales.

\section{Estructura Oficial del Documento de Análisis de Mercado}

Para que el informe final cumpla con los estándares académicos e
institucionales exigidos, se utilizará la siguiente tabla de contenidos:

\begin{enumerate}
    \item \textbf{Introducción y Justificación} \\
    Antecedentes del estudio de mercado, objetivos (general y específicos)
    y justificación frente a los criterios de pertinencia y acreditación
    ARCUSUR.

    \item \textbf{Marco Metodológico} \\
    Diseño cuantitativo y exploratorio, definición de la población
    objetivo, técnicas de recolección de datos y especificación del
    pipeline de procesamiento estadístico (Data Wrangling, estadística
    univariada, bivariada con pruebas $\chi^2$ y visualización
    computacional).

    \item \textbf{Análisis de la Demanda Educativa y Perfil de los
    Titulados} \\
    Caracterización demográfica, temporal (media y desviación estándar de
    la antigüedad de egreso) y presentación de matrices de datos e
    indicadores de formación continua y fuentes de financiamiento.

    \item \textbf{Análisis de la Oferta Laboral y Requerimientos del
    Sector Productivo} \\
    Mapeo de áreas tecnológicas de alta demanda, evaluación de
    competencias técnicas/transversales y ejemplos prácticos de
    modelamiento analítico.

    \item \textbf{Diagnóstico y Síntesis de Resultados} \\
    Cruces de variables clave, matrices de frecuencia cruzada y hallazgos
    principales para la Dirección de Carrera.

    \item \textbf{Conclusiones y Propuestas Estratégicas} \\
    Líneas de acción orientadas a los egresados (fomento de becas,
    posgrados modulares, vinculación empresarial B2B y observatorio
    permanente).

    \item \textbf{Referencias Bibliográficas} (Según Normativa APA 7\textsuperscript{a} edición).
\end{enumerate}

\section{Especificaciones Técnicas para el Desarrollo de la Herramienta de Recolección de Datos}

La herramienta de recolección (formulario web / sistema de encuestas) debe
construirse para capturar de forma limpia las variables requeridas en el
análisis. A continuación, se detallan las Historias de Usuario clave que
guiarán su desarrollo, junto con las directrices exactas de interpretación
analítica de la información generada.

\subsection*{Historia de Usuario 1: Registro de Antigüedad y Empleabilidad}

\textbf{Enunciado:} Como titulado de Ingeniería en Sistemas, quiero
registrar mi año de egreso, estado de ocupación actual y sector de
desempeño en la herramienta, para que la dirección de carrera conozca el
perfil temporal y la tasa de inserción de mi cohorte.

\textbf{Implementación técnica sugerida:} Campos obligatorios de tipo
numérico para el año de egreso y menús desplegables cerrados (dropdowns)
para sector y estado laboral, evitando sesgos por texto libre.

\textbf{Interpretación adecuada de la información:} Permite calcular
estadísticos descriptivos univariados (media, mediana y desviación
estándar de los años transcurridos desde la titulación). Su lectura indica
el grado de madurez profesional de la muestra, permitiendo segmentar a los
egresados entre profesionales junior (alta demanda de especializaciones
iniciales) y mandos consolidados.

\subsection*{Historia de Usuario 2: Preferencias de Formación Continua y Áreas de Interés}

\textbf{Enunciado:} Como titulado, quiero indicar el nivel de posgrado de
mi interés (Maestría, Especialidad, Doctorado) y el área temática
prioritaria (ej. Ciencia de Datos, Ciberseguridad, Desarrollo de
Software), para que la universidad oriente su oferta académica hacia las
tecnologías con mayor demanda real.

\textbf{Implementación técnica sugerida:} Preguntas de selección múltiple
con límite o selección única jerárquica para evitar ambigüedades en la
categorización de datos.

\textbf{Interpretación adecuada de la información:} Genera distribuciones
de frecuencia absoluta y relativa. Permite identificar concentraciones
críticas de mercado (por ejemplo, cuando ramas como Ciencia de Datos e
Ingeniería de Software agrupan más del 70\% de las preferencias), lo que
justifica técnica y financieramente la apertura de nuevos programas
académicos modulares o híbridos.

\subsection*{Historia de Usuario 3: Identificación de Mecanismos de Financiamiento}

\textbf{Enunciado:} Como titulado, quiero reportar las fuentes de
financiamiento estimadas para mis futuros estudios (recursos propios,
becas estatales, créditos educativos o patrocinio empresarial), para que
la institución detecte restricciones económicas y gestione alianzas
estratégicas.

\textbf{Implementación técnica sugerida:} Opción de selección categórica
nominal con categorías mutuamente excluyentes o ponderadas.

\textbf{Interpretación adecuada de la información:} Facilita la
construcción de tablas de contingencia cruzadas con el tipo de posgrado
deseado y la aplicación de pruebas de independencia (Chi-Cuadrada). Su
correcta interpretación revela si existe dependencia financiera (ej. alta
vulnerabilidad ligada a recursos propios o becas limitadas), orientando
directrices directas como la firma de convenios con banca y la
diversificación de becas institucionales.

\subsection*{Historia de Usuario 4: Evaluación de Brechas en Competencias Profesionales}

\textbf{Enunciado:} Como titulado con experiencia en el mercado, quiero
evaluar la suficiencia de las competencias técnicas (hard skills) y
transversales (soft skills) adquiridas en mi formación de grado frente a
los requerimientos de mi puesto, para retroalimentar la malla curricular.

\textbf{Implementación técnica sugerida:} Escalas Likert (de 1 a 5)
combinadas con campos de texto estructurado opcional.

\textbf{Interpretación adecuada de la información:} Proporciona un
diagnóstico cuantitativo del nivel de satisfacción y de las brechas
formativas detectadas en el mercado laboral. Estos puntajes medios y
desviaciones estándar actúan como evidencia empírica directa requerida por
los pares evaluadores en los procesos de acreditación de calidad.

\subsection*{Historia de Usuario 5: Automatización de Reportes y Matrices de Frecuencia Cruzada}

\textbf{Enunciado:} Como administrador del sistema y miembro del comité de
evaluación, quiero exportar matrices de frecuencia cruzada y gráficos
estadísticos interactivos (mapas de calor y diagramas de radar), para
sustentar con rigor matemático el Plan de Desarrollo Estratégico.

\textbf{Implementación técnica sugerida:} Módulo de exportación de datos
en formato tabular limpio (CSV/Excel) y un panel analítico (dashboard)
automatizado en el backend.

\textbf{Interpretación adecuada de la información:} Sintetiza la fase de
Data Wrangling y análisis exploratorio computacional. Traduce
interacciones de datos complejos en reportes listos para la toma de
decisiones gerenciales y académicas, garantizando trazabilidad y
transparencia en los compromisos de mejora continua de la carrera.

\section{Producto Final}

La guía metodológica y técnica resultará en un documento final de Análisis
de Mercado sobre el perfil, expectativas de formación y empleabilidad de
los titulados de la carrera de Ingeniería en Sistemas.

Así mismo, al final se deberá contar con una herramienta de software de
colección de datos que permita experimentar y analizar diferentes
escenarios que permitan entender y apoyar la toma de decisiones de gestión
académica que potencia permanentemente la capacidad académica y de fuerza
laboral de los profesionales que salen de la Universidad Mayor de San
Simón.

\section{Recomendaciones y Plazos}

Este enunciado es referencial por lo que el estudiante deberá respaldar
fehacientemente la toma de datos en campo en relación a las empresas e
industrias vinculadas de otra manera a la tecnología computacional. Esta
actividad requiere individualizar al menos tres empresas de desarrollo de
software o del ámbito dedicado al negocio de las tecnologías, una del
ámbito público y otra del ámbito privado.

Respecto de la fecha de entrega de esta actividad es el 13 de septiembre
de los corrientes, de acuerdo a la modalidad ya practicada en la actividad
anterior.

\section{Anexos}

La encuesta propuesta es la siguiente:

\subsection*{Encuesta de opinión a titulados de Ingeniería en Sistemas}

\textit{Estimad@ Ingenier@: La Carrera de Ingeniería en Sistemas está en
proceso de Autoevaluación para la elaboración de un PLAN DE DESARROLLO
ESTRATÉGICO en el marco de una futura postulación de acreditación ante el
ARCUSUR. Por lo que debido a la gran importancia de este proceso le
pedimos por favor que pueda ayudarnos compartiendo su opinión sobre varios
aspectos importantes de la Carrera.}

\subsubsection*{0. Datos generales}

\begin{itemize}[nosep]
    \item Edad que tiene actualmente (opción única, hasta ``35 años o más'').
    \item Género (opción única): Hombre / Mujer / Otro.
\end{itemize}

\subsubsection*{Formación complementaria}

\begin{itemize}[nosep]
    \item ¿Ha realizado o se encuentra realizando algún programa de
    formación complementaria (posgrado, diplomado, doble titulación,
    especialización, etc.)? SI / NO
    \item ¿Cuál es el programa de formación complementaria de mayor nivel
    que ha concluido o que se encuentra realizando? Diplomado /
    Especialidad / Maestría / Doctorado / Pos doctorado.
    \item ¿Dónde ha cursado el programa indicado? EUPG / Dirección de
    Posgrado FCyT / IBNORCA / Universidades internacionales /
    Universidades públicas nacionales / Universidades privadas nacionales
    / Otro.
    \item ¿Cuál fue la fuente de financiamiento para cursar el programa
    indicado? Financiado con recursos propios / Total financiamiento por
    medio de una beca / Financiado parcialmente por medio de una beca.
    \item Indique el nombre del programa (texto libre).
\end{itemize}

\subsubsection*{Interés en posgrado}

\begin{itemize}[nosep]
    \item ¿Estaría interesado en realizar estudios en algún programa de
    posgrado dentro de las áreas de la Ingeniería en Sistemas? SI / NO
    \item ¿En cuál de las siguientes áreas le interesaría realizar el
    programa? IA / Robótica / Ingeniería de Software / Base de Datos /
    Redes de Datos y Seguridad.
    \item En caso de realizar un estudio de posgrado ¿en qué organización
    educativa optaría por realizar sus estudios? (mismas opciones que
    ``¿Dónde ha cursado el programa indicado?'').
    \item Opinión (escala Likert: Totalmente en desacuerdo / En
    desacuerdo / De acuerdo / Totalmente de acuerdo) sobre: ¿Considera que
    existen programas de formación de posgrado ofertados en la Dirección
    de Posgrado de la FCyT que son afines a las áreas de especialidad de
    un Ingeniero en Sistemas?
\end{itemize}

\subsubsection*{Trayectoria profesional}

\begin{itemize}[nosep]
    \item ¿Cuántos son los años de vida profesional que lleva hasta la
    fecha? (numérico)
    \item De esa cantidad, ¿qué cantidad corresponde al total de tiempo
    en que se ha encontrado desempleado? (numérico)
    \item ¿Cuál de las siguientes afirmaciones describe su situación
    laboral actual? Actualmente trabajando en una organización /
    Actualmente NO trabajo / Dedico mi tiempo completo a un emprendimiento
    propio.
    \item ¿Cuál es el origen del desarrollo de su propio emprendimiento?
    Idea propia / Idea conjunta con compañeros / StartUp / Spin Off
    (derivado de un proyecto de la Universidad) / Otro.
    \item ¿Qué tipo de entregable genera el negocio? Producto / Servicio
    / Ambos.
    \item ¿Ha requerido financiamiento externo para su negocio? Sí,
    mediante préstamos bancarios / Sí, mediante inversores / No,
    financiamiento propio.
    \item Escala Likert (Insatisfecho / Algo satisfecho / Satisfecho):
    ¿Qué tan satisfecho está con el rendimiento actual de su negocio?
    \item Escala Likert (No es importante / Poco importante / Importante
    / Muy importante): ¿Qué tan importante considera que fue su formación
    en Ingeniería en Sistemas para iniciar y gestionar su negocio?
    \item ¿El trabajo que ejerce actualmente es su primer empleo? SI / NO
    \item Luego de titularse, ¿cuánto tiempo se demoró en conseguir su
    primer trabajo relacionado con su formación? Ya trabajaba antes de
    titularse / Menos de 1 mes / Entre 1--4 meses / Entre 4--8 meses /
    Entre 8--12 meses / Más de 12 meses.
    \item Desde que se tituló, ¿en cuántos empleos se ha desempeñado?
    (incluyendo el actual) 1 / 2 / 3 / Más de 3.
\end{itemize}

\subsubsection*{Satisfacción con la formación}

\begin{itemize}[nosep]
    \item Escala Likert (Totalmente en desacuerdo / En desacuerdo / De
    acuerdo / Totalmente de acuerdo): De manera global, ¿cuál es el grado
    de satisfacción que ha desarrollado con la formación profesional
    recibida en la Carrera de Ingeniería de Sistemas?
    \item Escala Likert (misma escala): Considera que existe concordancia
    entre la formación profesional recibida en la Carrera y los
    requerimientos (del mercado laboral).
    \item ¿Qué aspectos de la Carrera les resultaron de bastante utilidad?
    (elección múltiple): aplicación práctica de conceptos teóricos en
    proyectos reales / oportunidad de pasantías o prácticas profesionales
    / distribución de horarios de clases / desarrollo de habilidades
    analíticas y de resolución de problemas / colaboración de equipos
    multidisciplinarios / enseñanza de herramientas y software específico
    de la industria.
    \item ¿Qué aspectos de la Carrera considera que pueden mejorarse?
    (elección múltiple): mayor integración teoría-práctica / mayor
    relacionamiento con empresas por prácticas profesionales /
    distribución de horarios (horario para trabajadores, módulos
    virtuales, etc.) / actualización constante del plan de estudios /
    incremento de actividades extracurriculares (investigación, servicio
    social, etc.) / ofrecimiento de más opciones de especialización.
    \item Indique 5 asignaturas cursadas que considera que le brindaron
    una ventaja competitiva (elección múltiple, opciones: Introducción /
    Base de Datos / Inteligencia Artificial / Redes de Computadora /
    Electivas / T.I.S. / Básicas (Física, Cálculo, Álgebras) /
    Planificación de Proyectos y Gestión / Tecnología Redes Avanzadas /
    otra asignatura -- texto libre).
    \item Indique 5 asignaturas cursadas que no le resultaron de mucha
    utilidad (elección múltiple, opciones: Básicas (Física, Cálculo,
    Álgebras) / Aplic. Interactivas para Televisión Digital / Telefonía
    IP / Web Semánticas / Graficación por Computadora / Robótica / Taller
    de Programación en Bajo Nivel / Teoría de Autómatas y Leng. Formales /
    Programación Funcional / Teoría de Grafos / Contabilidad Básica /
    otra asignatura -- texto libre).
\end{itemize}

\subsubsection*{Red de contactos}

\begin{itemize}[nosep]
    \item ¿Le interesaría formar parte de una red de contactos oficial de
    la Carrera? SI / NO
    \item Email de contacto (texto libre).
    \item Número telefónico de contacto, de preferencia con WhatsApp
    (texto libre).
    \item Link de LinkedIn (texto libre).
\end{itemize}

\subsection*{Encuesta a empleadores}

\textit{La Carrera de Ingeniería de Sistemas de la Facultad de Ciencias y
Tecnología (UMSS), en el marco de su proceso de Autoevaluación Académica,
desarrolla la etapa de recopilación de información primaria mediante la
aplicación del presente cuestionario dirigido a empleadores de
profesionales titulados de la Carrera. El objetivo es conocer la
percepción del sector empleador respecto al desempeño profesional,
competencias, capacidades y formación académica de los Ingenieros de
Sistemas titulados por la UMSS, con la finalidad de identificar
fortalezas y oportunidades de mejora en el proceso formativo de la
Carrera. La información será utilizada exclusivamente con fines
académicos e institucionales, garantizando confidencialidad y anonimato.}

\begin{enumerate}
    \item Correo (obligatorio).
    \item Nombre de la Empresa o Institución (obligatorio, texto libre).
    \item Tipo de organización (obligatorio): Público / Privado.
    \item Tamaño de la organización (obligatorio): Grande (100+
    funcionarios) / Mediana (31--99) / Pequeña (11--30) / Micro (hasta 10).
    \item Rubro o sector principal de la organización (obligatorio,
    opción única): Desarrollo de software / Servicios de TI /
    Telecomunicaciones y redes / Soporte técnico y mantenimiento de
    sistemas / Consultoría en sistemas / transformación digital / Banca y
    servicios financieros / Instituciones públicas / gobierno / Educación
    e investigación / Industria y manufactura / Comercio y comercio
    electrónico / Salud / Seguridad informática / ciberseguridad /
    Análisis de datos / inteligencia de negocios / Startups o
    emprendimientos tecnológicos / Otros.
\end{enumerate}

\subsubsection*{Empleabilidad}

Se busca conocer la demanda laboral, oportunidades de contratación y
desempeño de los profesionales titulados de Ingeniería de Sistemas en las
organizaciones, con el fin de evaluar la pertinencia de la formación
profesional frente a las necesidades del entorno laboral.

\begin{enumerate}[resume]
    \item ¿La organización ha contratado Ingenieros de la carrera de
    Sistemas en los últimos 5 años? SI / NO
    \item Considerando las necesidades actuales y proyectadas, ¿existe la
    posibilidad de incorporar Ingenieros de Sistemas en los próximos
    años? Alta probabilidad de contratación / Depende del crecimiento o
    nuevos proyectos / No se tiene prevista la contratación actualmente /
    No requiere profesionales del área.
    \item ¿Qué nivel de formación académica o actualización profesional
    demanda su organización para contratar nuevos profesionales del
    área? Licenciatura / Diplomado / Especialidad / Maestría / Doctorado
    / Post-doctorado.
    \item ¿A través de qué medio convoca a profesionales en su
    organización? Medios escritos y audiovisuales / Internet y redes
    sociales / Postulación competitiva (convocatorias públicas y
    privadas) / Oferta profesional personal / Recomendación de terceros /
    Otros.
    \item ¿Qué tipo de cargos desempeñan los profesionales titulados de
    la Carrera en su organización? (selección múltiple, obligatorio): de
    apoyo / asistencia tecnológica / técnico (desarrollo de software,
    administración de sistemas, redes, bases de datos, soporte TI) /
    analista (sistemas, datos, procesos o soluciones tecnológicas) /
    especialista (ciberseguridad, arquitectura de software,
    infraestructura tecnológica, etc.) / supervisor / coordinador de área
    tecnológica / jefatura o gerencial (gestión de proyectos, dirección
    de TI, transformación digital) / consultor tecnológico / otro.
    \item Desde su experiencia en el sector laboral, ¿qué nuevas áreas de
    conocimiento, tecnologías o tendencias emergentes considera
    importantes para la formación del profesional? (texto libre).
    \item ¿Cuáles son las nuevas herramientas tecnológicas (software,
    plataformas o entornos de desarrollo) que debería conocer y manejar
    el profesional titulado? (texto libre).
    \item Desde su experiencia, ¿qué habilidades, competencias y
    destrezas considera fundamentales que debe desarrollar el profesional
    de Ingeniería de Sistemas? (texto libre).
\end{enumerate}

\subsubsection*{Formación y perfil profesional}

Se busca conocer la valoración de los empleadores sobre la formación,
competencias y desempeño profesional de los titulados de Ingeniería de
Sistemas, con el propósito de fortalecer la calidad académica y adecuar
el perfil profesional a las necesidades del mercado laboral. Escala:
Totalmente de acuerdo / Parcialmente de acuerdo / Totalmente en
desacuerdo / No sabe. Criterios evaluados:

\begin{itemize}[nosep]
    \item La Carrera de Ingeniería de Sistemas de la UMSS brinda
    confianza a nuestra organización como formadora de profesionales del
    área tecnológica.
    \item Considero que el título profesional otorgado por la Carrera es
    consistente con el nivel de formación y las competencias demostradas
    por sus titulados.
    \item Las autoridades de la Carrera consultan regularmente nuestra
    opinión como empleadores respecto al perfil profesional y las
    competencias requeridas en el ámbito laboral.
    \item Conozco el perfil profesional y las competencias de egreso
    declaradas por la Carrera.
    \item El perfil profesional declarado por la Carrera es coherente con
    el desempeño laboral demostrado por sus titulados.
    \item La formación proporcionada por la Carrera permite a sus
    titulados alcanzar un desempeño destacado tanto en el ámbito
    profesional como en su formación integral.
    \item La Carrera considera nuestra opinión como empleadores e
    incorpora las necesidades del sector laboral en la definición y
    mejora de actividades de formación práctica, pasantías y modalidades
    de titulación.
    \item El desempeño profesional de los titulados evidencia el
    desarrollo de competencias laborales como: comunicación efectiva,
    análisis y resolución de problemas, pensamiento crítico, aprendizaje
    continuo, trabajo en equipo, iniciativa personal y manejo adecuado de
    tecnologías de información.
    \item Los profesionales titulados demuestran valores y actitudes
    relacionados con la ética profesional, responsabilidad social,
    inclusión, diversidad, respeto a los derechos humanos y compromiso
    con el medio ambiente.
    \item He participado en procesos de recopilación de información,
    evaluación o retroalimentación relacionados con el desempeño laboral
    de los titulados de la Carrera.
\end{itemize}

\end{document}
    